\documentclass[12pt]{article}
\usepackage[T1]{fontenc}
\usepackage[utf8]{inputenc}
\usepackage{lmodern}
\usepackage{textcomp}
\usepackage{enumitem}
\usepackage{geometry}
\geometry{margin=1in}
\begin{document}
\begin{center}
Answer Key - Sociology - Graduate Level Assessment 21 \
\small (Answers are ordered exactly as in the question paper.)
\end{center}

\section*{Multiple Choice}
\begin{enumerate}[label=\arabic*.]
\item \textbf{Answer:} C
\\ \textit{Options:} A) cults of the capital; B) capital culture; C) cultural capital; D) culpable capture
\\ \textit{Note:} Bourdieu argued that cultural capital, which encompasses non-financial social assets such as education, intellect, and cultural knowledge, plays a crucial role in maintaining and reproducing class distinctions by influencing individuals' social mobility and access to resources.
\item \textbf{Answer:} C
\\ \textit{Options:} A) absolute poverty; B) relative poverty; C) vertical mobility; D) horizontal mobility
\\ \textit{Note:} Vertical mobility refers to the movement up or down the social hierarchy, indicating a change in social position, which aligns with the scenario of someone gaining a higher or lower social status than their birth position.
\item \textbf{Answer:} A
\\ \textit{Options:} A) triangulation; B) explanation; C) description; D) exploration
\\ \textit{Note:} Triangulation is a methodological approach used to enhance the validity and reliability of research findings, rather than a distinct purpose of research itself.
\item \textbf{Answer:} D
\\ \textit{Options:} A) a reduction in population size, caused by a higher rate of emigration than immigration; B) a change in the principal causes of death and disease since industrialization; C) increased birth and death rates, resulting in a relatively young population; D) a decline in the birth rate, greater life expectancy, and an ageing population
\\ \textit{Note:} The 'demographic transition' describes the shift in a society from high birth and death rates to lower birth rates, increased life expectancy, and an ageing population as it develops economically and socially.
\item \textbf{Answer:} B
\\ \textit{Options:} A) repress it, by promoting only the interests of elite groups; B) revive it, by reaffirming a commitment to freedom of speech; C) reproduce it, by emphasizing face-to-face contact with peer groups; D) replace it with a superior form of communication
\\ \textit{Note:} The correct answer is justified because the Internet has facilitated broader access to information and diverse viewpoints, thereby enhancing public discourse and reinforcing the principle of freedom of speech within the public sphere.
\end{enumerate}

\section*{True / False}
\begin{enumerate}[label=\arabic*.]
\setcounter{enumi}{5}
\item \textbf{Answer:} True
\\ \textit{Options:} A) True; B) False
\\ \textit{Note:} The statement is true because the emergence of the middle classes in the nineteenth century was characterized by urbanization, increased participation in civic life, and the prevalence of white-collar occupations, which collectively defined their social and economic status in increasingly industrialized societies.
\item \textbf{Answer:} False
\\ \textit{Options:} A) True; B) False
\\ \textit{Note:} Ethnic identity is shaped not only by governmental classifications but also by personal experiences, cultural practices, and social interactions, making it a complex and subjective construct rather than a solely objective one.
\item \textbf{Answer:} True
\\ \textit{Options:} A) True; B) False
\\ \textit{Note:} Policy makers believed that former slaves migrating to urban areas would assimilate into mainstream society due to the expectation that urban environments would provide greater economic opportunities and social integration, thereby facilitating their integration into the dominant culture.
\item \textbf{Answer:} True
\\ \textit{Options:} A) True; B) False
\\ \textit{Note:} Laqueur's 'two-sex' model asserts that male and female are understood as fundamentally different and separate biological categories, emphasizing the distinctness of the sexes in terms of their roles and characteristics within society.
\item \textbf{Answer:} False
\\ \textit{Options:} A) True; B) False
\\ \textit{Note:} In the traditional hierarchy, the King's Grandsons hold a higher rank than the Archbishop of York and the Lord High Treasurer due to their direct royal lineage and connection to the monarchy.
\end{enumerate}

\section*{Long Form Response}
\begin{enumerate}[label=\arabic*.]
\setcounter{enumi}{10}
\item \textbf{Answer:} Townsend's (1979) indicators of relative deprivation highlight the disparities in access to resources and opportunities that contribute to social inequalities. One such indicator is the purchasing of consumer goods, exemplified by the act of buying fewer than twenty DVDs in the previous year. This behavior reflects not only an individual's financial constraints but also their exclusion from cultural participation compared to peers who can afford more. The limited ability to engage in leisure activities, such as watching films, illustrates broader systemic inequalities, where socioeconomic status dictates lifestyle choices and access to cultural capital. Such indicators reveal the multifaceted nature of deprivation, emphasizing that social inequalities extend beyond mere income disparities to include the quality of life and social integration.
\\ \textit{Note:} correct because it illustrates how Townsend's indicators of relative deprivation, such as limited consumer purchasing, serve as a measure of social inequalities by highlighting individuals' financial constraints and their resulting exclusion from cultural and leisure activities compared to more affluent peers.
\item \textbf{Answer:} National surveys play a crucial role in understanding societal trends in the UK by providing empirical data that reflect public attitudes, behaviors, and demographic changes. These surveys, such as the British Social Attitudes Survey and the Labour Force Survey, offer insights into various aspects of life, including economic conditions, health, and social cohesion. In contrast, the Fashion Sensibility Survey, despite its potential value, is not regularly conducted by the government, likely due to its niche focus and the perception that fashion preferences are more subjective and less critical to policy-making compared to broader societal issues. This absence indicates a prioritization of surveys that yield data directly applicable to social welfare and economic stability, illustrating the government's strategic emphasis on pressing social matters over lifestyle trends.
\\ \textit{Note:} correct because it emphasizes the importance of national surveys in providing empirical data on societal trends, while also noting that the Fashion Sensibility Survey, despite its potential insights, is not prioritized by the government, likely due to its niche focus compared to broader, more impactful surveys.
\end{enumerate}

\vfill
\noindent\textit{End of Answer Key}
\end{document}